\begin{solapartat}
La segona llei de Newton estableix que: $\vec{F} = M_{Ll}\vec{a}$ \quad \punts{0,1}

Segons la llei de gravitació universal, el mòdul de la força s'expressa com:

$F = G\,\dfrac{M_{Ll}M_T}{r^2}$ \quad \punts{0,2}

Per tant obtenim que: $a = G\,\dfrac{M_T}{r^2}$ \quad \punts{0,1}

D'altra banda, considerant que el satèl·lit descriu un moviment circular uniforme al voltant de al terra, la seva acceleració centrípeta és: $a = \dfrac{v^2}{r}$ o $a = \omega^2 r$ \punts{0,1}, i la velocitat es pot expressar com $v = \dfrac{2\pi r}{T}$ o $\omega = \dfrac{2\pi}{T}$ \quad \punts{0,1}

Per tant: $T = 2\pi\sqrt{\dfrac{r^3}{GM_T}}$ \quad \punts{0,1}

Llavors: $T = 2\pi\sqrt{\dfrac{(3,84\times 10^{8})^3}{6,67\times 10^{-11}\cdot 5,98\times 10^{24}}} = 2,37\times 10^{6}\ \mathrm{s} = 27,4\ \text{dies}$ \quad \punts{0,1}

\punts{0,35} $U_g = -G\,\dfrac{M_{Ll}M_T}{r} = -6,67\times 10^{-11}\,\dfrac{7,35\times 10^{22}\cdot 5,98\times 10^{24}}{3,84\times 10^{8}} = -7,63\times 10^{28}\ \mathrm{J}$ \quad \punts{0,1}
\end{solapartat}

\begin{solapartat}
Segons la llei de gravitació universal,

\punts{0,4} $G\,\dfrac{mM_T}{d^2} = G\,\dfrac{mM_{Ll}}{(R_{TLl}-d)^2}$

\punts{0,5} $\left[\dfrac{R_{TLl}-d}{d}\right]^2 = \dfrac{M_{Ll}}{M_T} \Longrightarrow \dfrac{R_{TLl}-d}{d} = \pm\sqrt{\dfrac{M_{Ll}}{M_T}} \Longrightarrow \dfrac{R_{TLl}}{d} = 1 \pm \sqrt{\dfrac{M_{Ll}}{M_T}} \Longrightarrow d = \dfrac{R_{TLl}}{1 \pm \sqrt{\frac{M_{Ll}}{M_T}}}$

Dues possibles solucions:

% DUBTE: amb les dades, R_TLl/(1 - sqrt(M_Ll/M_T)) = 4,32e8 m (la pauta diu 4,31e8 m).
\punts{0,35} $d = \dfrac{R_{TLl}}{1-\sqrt{\frac{M_{Ll}}{M_T}}} = 4,31\times 10^{8}\ \mathrm{m} > R_{TLl}$, aquesta solució no es bona, atès que no es troba entre la Terra i la Lluna tal i com demana l'enunciat.

$d = \dfrac{R_{TLl}}{1+\sqrt{\frac{M_{Ll}}{M_T}}} = 3,46\times 10^{8}\ \mathrm{m} < R_{TLl}$, aquesta és la resposta correcte!
\end{solapartat}
